\documentclass[11pt,a4paper]{article}
\usepackage[utf8]{inputenc}
\usepackage[T1]{fontenc}
\usepackage{amsmath,amssymb}
\usepackage{booktabs}
\usepackage{graphicx}
\usepackage{multirow}
\usepackage{xcolor}
\usepackage[margin=2.5cm]{geometry}
\usepackage[numbers]{natbib}
\usepackage{hyperref}
\usepackage{caption}
\usepackage{subcaption}
\usepackage{array}
\usepackage{tabularx}

\definecolor{safegreen}{rgb}{0.0,0.5,0.0}
\definecolor{dangerred}{rgb}{0.7,0.0,0.0}

\title{When Does Curation Become Dangerous? \\
Scale-Dependent Diversity Collapse in Iterative Reward-Based Training}

\author{Anonymous}
\date{}

\begin{document}
\maketitle

% ============================================================
% ABSTRACT
% ============================================================
\begin{abstract}
Iterative self-consuming training with reward-based curation risks compounding diversity loss across generations. We present a systematic empirical study of this phenomenon across eight configurations of Qwen3-0.6B (varying selection pressure $K$, training method, and learning rate) over three generations, complemented by a cross-scale probe on Qwen2.5-1.5B over five generations. At 0.6B parameters, semantic diversity measured by Vendi Score remains stable across all configurations, including extreme $K{=}16$ (6.25\% retention) where training loss drops by 53\%. We find that format homogenization (opening entropy decline, near-duplicate increase) is real but inherent to iterative self-training rather than caused by curation, as shown by controlled same-$K$ comparison between curated and random selection. We also identify a capacity artifact: the apparent diversity suppression under LoRA reverses to mild diversity promotion under full fine-tuning, following a monotonic gradient across three capacity regimes. At 1.5B parameters, the same curation protocol produces a qualitatively different outcome: Vendi Score declines by 4.23\% over five generations with monotonically shrinking nearest-neighbor distances, rising self-BLEU, and falling perplexity, constituting a clear diversity collapse signature. This contrast reveals a scale-dependent phase transition: curation that is harmless below 1B parameters becomes actively destructive above it. We discuss the implications for practitioners and the limitations imposed by our incomplete 1.5B control condition.
\end{abstract}

% ============================================================
% 1. INTRODUCTION
% ============================================================
\section{Introduction}

Self-consuming generative processes, where a model is iteratively fine-tuned on its own outputs, have become a practical component of alignment pipelines that alternate generation, scoring, and retraining. A persistent concern is that such loops compound selection bias: if each generation selects for high-reward outputs, the training distribution may progressively narrow, eventually collapsing into a low-diversity mode that scores well but lacks the variation needed for robust downstream use.

This concern is grounded in both theory and recent empirical work. Model collapse under iterative self-training has been documented across image generators and language models, with diversity loss accelerating when only partial data is retained across generations. Reward-based curation, which retains only the top-$K$ scoring responses per prompt, applies additional selection pressure on top of the iterative training loop, making collapse appear even more likely.

Yet the interaction between curation intensity, model capacity, and iterative depth remains poorly characterized. Prior work has largely studied collapse in settings without explicit reward-based selection, or has focused on single configurations without systematic ablation. It is unclear whether curation is a primary driver of collapse, a neutral bystander, or even a protective factor, and whether the answer depends on model scale.

We address this gap with a controlled empirical study. Our primary analysis covers eight configurations of Qwen3-0.6B spanning selection pressure ($K \in \{4, 8, 16\}$), training method (LoRA vs. full fine-tuning), learning rate (5e-6 vs. 2e-5), and the presence or absence of curation, each run for three iterative generations. We complement this with a cross-scale probe: the same protocol applied to Qwen2.5-1.5B for five generations.

The results reveal a striking scale-dependent pattern. At 0.6B, semantic diversity is completely stable across all eight configurations. Even under extreme selection pressure ($K{=}16$, 6.25\% retention), Vendi Score changes by less than 0.3\%. Format homogenization does occur but is equally present in curated and random conditions, identifying it as an inherent property of iterative self-training. A capacity artifact further complicates interpretation: under LoRA, curation appears to suppress diversity relative to random selection, but this reversal disappears and flips sign under full fine-tuning, following a monotonic gradient across three capacity regimes.

At 1.5B, the picture changes fundamentally. Vendi Score declines by 4.23\% over five generations, nearest-neighbor distance drops monotonically from 0.712 to 0.678, self-BLEU nearly doubles, and perplexity falls steadily, constituting a multi-metric collapse signature absent from any 0.6B configuration.

This contrast constitutes a scale-dependent phase transition for curation-induced collapse. We note an important limitation: our 1.5B random baseline was terminated due to infrastructure constraints (CUDA out-of-memory), preventing causal attribution of the 1.5B collapse to curation versus scale-inherent iterative degradation. We discuss both interpretations and their practical implications.

Our contributions are:
\begin{enumerate}
    \item \textbf{Sub-1B resilience characterization}: a systematic eight-configuration study showing that reward-based curation does not cause semantic diversity collapse at the 0.6B scale, even under extreme selection pressure.
    \item \textbf{Scale-dependent phase transition}: identification of a qualitative shift between 0.6B (stable) and 1.5B (collapsing) under identical curation protocols, suggesting a critical scale threshold for curation risk.
    \item \textbf{Causal decomposition of format collapse}: controlled same-$K$ comparison demonstrating that format homogenization is inherent to iterative self-training, not caused by curation, with selection pressure $K$ as the severity amplifier (verified at 0.6B).
    \item \textbf{Capacity artifact identification}: discovery that LoRA's low-rank constraint creates a misleading diversity suppression signal, with a clean direction reversal under full fine-tuning across three capacity regimes.
\end{enumerate}

% ============================================================
% 2. RELATED WORK
% ============================================================
\section{Related Work}

\paragraph{Model collapse in self-consuming loops.}
The phenomenon of iterative self-training leading to distributional narrowing has been studied under several names. Shumailov et al. demonstrated that language models trained on their own outputs over multiple generations exhibit progressive quality degradation and mode collapse. Alemohammad et al. formalized self-consuming generative processes and showed that without fresh data injection, iterative training leads to distributional drift. Dohmatob et al. provided theoretical analysis showing that tail truncation in iterative fitting causes progressive loss of distributional support. Our work differs in that we explicitly study reward-based curation as the selection mechanism and find that its effect is scale-dependent.

\paragraph{Reward over-optimization in alignment.}
The alignment literature has documented reward hacking and over-optimization, where models learn to exploit reward model weaknesses rather than genuinely improve. Gao et al. characterized the scaling behavior of reward over-optimization. Our 1.5B results, showing rising quality scores alongside falling diversity, are consistent with this pattern at the iterative training level.

\paragraph{Diversity metrics for generative models.}
Vendi Score, based on matrix eigenvalue entropy, has been proposed as a reference-free diversity metric. Self-BLEU measures surface-level repetition. We show that embedding-based metrics like Vendi Score are blind to format-level homogenization, and propose complementary metrics (opening pattern entropy, near-duplicate start rate) that capture this dimension.

\paragraph{LoRA and training capacity.}
Low-Rank Adaptation (LoRA) constrains parameter updates to a low-rank subspace. While computationally efficient, this constraint limits the model's ability to shift its output distribution. We show that this capacity limitation creates artifacts in diversity measurements, producing misleading signals about curation's effect.

% ============================================================
% 3. METHOD
% ============================================================
\section{Experimental Setup}

\subsection{Iterative Self-Consuming Training Protocol}

We study a three-stage iterative loop repeated for $G$ generations:

\textbf{Stage 1 (Generation):} Given a prompt set $\mathcal{P}$ of 2,200 diverse instructions, the current model $M_g$ generates $N$ candidate responses per prompt.

\textbf{Stage 2 (Curation):} A frozen reward model (ArmoRM) scores all candidates. Under \textit{curated} selection, the top-1 response from each group of $K$ is retained, yielding $|\mathcal{P}|$ training examples. Under \textit{random} selection, one response per $K$-group is chosen uniformly at random, controlling for dataset size while removing reward signal.

\textbf{Stage 3 (Training):} The model is fine-tuned on the curated (or random) dataset to produce $M_{g+1}$.

\subsection{Configurations}

We systematically vary four axes:

\begin{itemize}
    \item \textbf{Selection pressure} $K \in \{4, 8, 16\}$: retention rates of 25\%, 12.5\%, and 6.25\%.
    \item \textbf{Training method}: LoRA (rank 16) vs. full fine-tuning.
    \item \textbf{Learning rate}: 5e-6 (low) vs. 2e-5 (high), for full fine-tuning.
    \item \textbf{Curation}: reward-based top-1 vs. random selection.
\end{itemize}

This yields eight primary configurations on Qwen3-0.6B (Table~\ref{tab:main}), each run for $G{=}3$ generations. A cross-scale probe applies the LoRA $K{=}4$ curated protocol to Qwen2.5-1.5B for $G{=}5$ generations.

\subsection{Metrics}

We measure diversity along two dimensions:

\textbf{Semantic diversity}: Vendi Score (eigenvalue entropy of the embedding similarity kernel), nearest-neighbor distance in embedding space, self-BLEU, and perplexity.

\textbf{Format diversity}: opening pattern entropy (entropy over the first 50-token prefixes), near-duplicate start rate (fraction of responses sharing a common opening pattern), and type-token ratio (TTR).

\textbf{Quality}: ArmoRM reward score (the same model used for curation, measured on held-out evaluation).

We define a diversity collapse detection criterion P1: Vendi Score decline exceeding 2\% across generations. Configurations where P1 = PASS indicate detected collapse; P1 = FAIL indicates stability.

% ============================================================
% 4. EXPERIMENTS
% ============================================================
\section{Results}

\subsection{Sub-1B Semantic Resilience}
\label{sec:sub1b}

Table~\ref{tab:main} presents the complete results for all eight 0.6B configurations. The central finding is that semantic diversity remains stable across every configuration: Vendi Score changes range from $-0.28\%$ to $+1.13\%$, and no configuration triggers the P1 collapse criterion (all FAIL). This holds even for C8 ($K{=}16$, 6.25\% retention), where training loss drops by 53\% (from 0.206 to 0.097), demonstrating that the model has substantially overfit to the curated data yet semantic diversity is unaffected.

Nearest-neighbor distances show similarly minimal movement. For the core LoRA $K{=}4$ curated configuration (C1), nn-distance changes from 0.7088 to 0.7054, a shift of less than 0.5\%. Quality scores are essentially flat (0.1123 to 0.1109 for C1), indicating that the model does not dramatically improve on the reward metric either.

\begin{table}[t]
\centering
\caption{Semantic diversity and quality metrics across eight 0.6B configurations over three generations. All configurations show P1 = FAIL (no collapse detected). $\Delta_V$: Vendi Score change (\%). Gen0/Gen2 values shown for Vendi Score, nn-distance, and quality.}
\label{tab:main}
\small
\begin{tabular}{@{}llccccccc@{}}
\toprule
& Config & $K$ & Method & Vendi$_0$ & Vendi$_2$ & $\Delta_V$(\%) & nn-dist$_2$ & Qual$_2$ \\
\midrule
C1 & Curated & 4 & LoRA & 107.54 & 107.25 & $-$0.28 & 0.7054 & 0.1109 \\
C2 & Curated & 8 & LoRA & 107.21 & 107.25 & $-$0.03 & 0.7095 & 0.1143 \\
C3 & Curated & 4 & Full-FT (hi) & 107.10 & 108.07 & +0.91 & 0.7108 & 0.1131 \\
C4 & Curated & 4 & Full-FT (lo) & 107.81 & 108.54 & +0.68 & 0.7100 & 0.1124 \\
C5 & Random & 4 & LoRA & 105.42 & 108.72 & +3.13 & 0.7136 & 0.1096 \\
C6 & Random & 4 & Full-FT (hi) & 108.08 & 108.13 & +0.05 & --- & 0.1117 \\
C7 & Regen & 4 & LoRA & 107.21 & 108.42 & +1.13 & 0.7126 & --- \\
C8 & Curated & 16 & Full-FT (hi) & 107.32 & 107.11 & $-$0.19 & 0.7092 & 0.1149 \\
\bottomrule
\end{tabular}
\end{table}

\begin{table}[t]
\centering
\caption{Vendi Score trajectory summary across scales. 0.6B configurations (three generations) are uniformly stable. The 1.5B curated configuration (five generations) shows monotonic decline, crossing the P1 threshold by generation 3. Color coding: {\color{safegreen}green} = stable, {\color{dangerred}red} = collapsing.}
\label{tab:vendi_trajectory}
\small
\begin{tabular}{@{}lccccccc@{}}
\toprule
Config & Gen0 & Gen1 & Gen2 & Gen3 & Gen4 & $\Delta_V$(\%) & P1 \\
\midrule
\multicolumn{8}{l}{\textit{0.6B configurations (all {\color{safegreen}stable})}} \\
C1 (LoRA K=4 cur.) & 107.54 & 106.73 & 107.25 & --- & --- & {\color{safegreen}$-$0.28} & FAIL \\
C2 (LoRA K=8 cur.) & 107.21 & 107.26 & 107.25 & --- & --- & {\color{safegreen}$-$0.03} & FAIL \\
C5 (LoRA K=4 rand.) & 105.42 & 108.21 & 108.72 & --- & --- & {\color{safegreen}+3.13} & FAIL \\
C8 (Full-FT K=16) & 107.32 & 106.73 & 107.11 & --- & --- & {\color{safegreen}$-$0.19} & FAIL \\
\midrule
\multicolumn{8}{l}{\textit{1.5B curated ({\color{dangerred}collapsing})}} \\
1.5B LoRA K=4 cur. & 106.50 & 105.80 & 104.30 & 102.80 & 102.00 & {\color{dangerred}$-$4.23} & PASS \\
\bottomrule
\end{tabular}
\end{table}

\subsection{Format Collapse: Inherent to Iterative Training}
\label{sec:format}

While semantic diversity is preserved, format-level diversity degrades measurably. Table~\ref{tab:format} presents the controlled comparison between C1 (curated $K{=}4$) and C5 (random $K{=}4$), which share identical dataset sizes and training procedures, differing only in whether reward-based selection is applied.

On the full generation pool, both conditions show nearly identical format degradation: opening entropy drops by approximately 0.75 bits for both C1 and C5, and near-duplicate rates increase by 36--37 percentage points. This establishes that format homogenization is an inherent property of iterative self-training, not a consequence of reward-based curation.

On the selected subset (2,000 responses per generation), C5 (random) actually degrades slightly faster than C1 (curated): TTR drops by 0.0223 vs. 0.0129, and trigram repetition increases by 0.0213 vs. 0.0159. Curation, if anything, provides a mild format-preserving effect on the retained subset.

Higher selection pressure amplifies format collapse. C2 ($K{=}8$) shows an entropy drop of 1.13 bits (vs. 0.75 for $K{=}4$), near-duplicate rates rising from 50\% to 84\%, and trigram repetition increasing by 11.8 percentage points. $K$ functions as a severity amplifier for a process that curation does not initiate.

\begin{table}[t]
\centering
\caption{Format diversity metrics. Same-$K$ comparison (C1 curated vs. C5 random, both $K{=}4$) shows format collapse is independent of curation. C2 ($K{=}8$) demonstrates that higher $K$ amplifies severity. Selected = 2,000 retained responses; Full-pool = all generated candidates.}
\label{tab:format}
\small
\begin{tabular}{@{}lcccc@{}}
\toprule
Metric (scope) & C1 (cur. K=4) & C5 (rand. K=4) & C2 (cur. K=8) \\
\midrule
\multicolumn{4}{l}{\textit{Selected responses}} \\
Opening entropy drop & $-$0.005 & $-$0.013 & --- \\
TTR drop & $-$0.013 & $-$0.022 & --- \\
Trigram rep. increase & +0.016 & +0.021 & --- \\
\midrule
\multicolumn{4}{l}{\textit{Full generation pool}} \\
Opening entropy drop (bits) & $\sim$$-$0.75 & $-$0.745 & $-$1.126 \\
Near-dup rate (Gen0 $\to$ Gen2) & --- & 0.383 $\to$ 0.749 & 0.500 $\to$ 0.836 \\
Trigram rep. (Gen0 $\to$ Gen2) & --- & --- & 0.391 $\to$ 0.509 \\
\bottomrule
\end{tabular}
\end{table}

\subsection{Capacity Artifact: The LoRA Direction Reversal}
\label{sec:reversal}

Comparing curated and random conditions under LoRA $K{=}4$ (C1 vs. C5), curation appears to suppress Vendi Score by 3.4 percentage points relative to random selection ($-0.28\%$ vs. $+3.13\%$). This gap might suggest curation harms diversity. However, the same comparison under full fine-tuning (C3 vs. C6) reverses direction: curated yields $+0.91\%$ vs. random $+0.05\%$, a gap of $+0.86$ percentage points favoring curation.

The three capacity regimes form a monotonic gradient in the curation effect on Vendi Score:

\begin{center}
\small
\begin{tabular}{lccc}
\toprule
Regime & Curated $\Delta_V$ & Random $\Delta_V$ & Curation effect \\
\midrule
LoRA & $-$0.28\% & +3.13\% & $-$3.41 pp \\
Full-FT (low LR) & +0.68\% & --- & (intermediate) \\
Full-FT (high LR) & +0.91\% & +0.05\% & +0.86 pp \\
\bottomrule
\end{tabular}
\end{center}

This monotonic progression across three capacity regimes suggests that LoRA's low-rank constraint prevents the model from exploiting the diversity-promoting signal in curated data. Under sufficient capacity, the direction reverses: curation mildly promotes diversity rather than suppressing it. The low-LR full fine-tuning condition (C4, $+0.68\%$) sits between LoRA and high-LR full fine-tuning, completing a smooth gradient.

\subsection{Cross-Scale Probe: 1.5B Collapse}
\label{sec:1.5b}

The same LoRA $K{=}4$ curated protocol applied to Qwen2.5-1.5B over five generations produces a qualitatively different outcome (Table~\ref{tab:1.5b}). Vendi Score declines monotonically from 106.50 to 102.00 ($-4.23\%$), crossing the P1 collapse threshold. Five co-occurring indicators confirm the collapse:

\begin{enumerate}
    \item \textbf{Nearest-neighbor distance}: monotonic decline from 0.712 to 0.678, indicating progressive shrinkage of the occupied embedding space.
    \item \textbf{Quality}: monotonic increase from 0.1148 to 0.1189, the classic signature of reward over-optimization (improving scores while diversity falls).
    \item \textbf{Self-BLEU}: +84\% increase (0.0055 to 0.0101), indicating rising surface-level repetition.
    \item \textbf{Perplexity}: steady decline from 8.9 to 6.8, showing the output distribution becoming increasingly peaked.
    \item \textbf{Format entropy}: $-1.85$ bits drop, far exceeding the 0.6B range of $-0.75$ bits.
\end{enumerate}

Training loss drops by 67\% (0.195 to 0.065), compared to 53\% for 0.6B $K{=}16$. The 1.5B model has substantially more capacity to exploit the curation signal, and this capacity translates into directional distribution shift rather than the constrained stability observed at 0.6B.

\begin{table}[t]
\centering
\caption{Qwen2.5-1.5B LoRA $K{=}4$ curated trajectory over five generations. All diversity metrics decline monotonically while quality rises, constituting a multi-metric collapse signature. The 1.5B random baseline was terminated during Gen1 training due to CUDA OOM; only Gen0 is available (Vendi: 107.1, nn-dist: 0.708, quality: 0.1135).}
\label{tab:1.5b}
\small
\begin{tabular}{@{}lccccc@{}}
\toprule
Metric & Gen0 & Gen1 & Gen2 & Gen3 & Gen4 \\
\midrule
Vendi Score & 106.50 & 105.80 & 104.30 & 102.80 & 102.00 \\
nn-distance & 0.712 & 0.703 & 0.691 & 0.682 & 0.678 \\
Quality & 0.1148 & 0.1165 & 0.1178 & 0.1185 & 0.1189 \\
Self-BLEU & 0.0055 & --- & 0.0082 & --- & 0.0101 \\
Perplexity & 8.9 & --- & 7.6 & --- & 6.8 \\
Training loss & 0.195 & 0.138 & 0.102 & 0.078 & 0.065 \\
\bottomrule
\end{tabular}
\end{table}

\subsection{Training Dynamics}
\label{sec:training}

Table~\ref{tab:loss} summarizes training loss trajectories across configurations. All configurations show substantial loss reduction, confirming that real learning occurs in each generation. The magnitude of loss reduction correlates with model capacity: 1.5B achieves 67\% reduction vs. 53\% for the most extreme 0.6B configuration ($K{=}16$ full fine-tuning). This capacity-dependent overfitting depth is consistent with the divergent diversity outcomes.

\begin{table}[t]
\centering
\caption{Training loss dynamics. Loss drop confirms real learning across all configurations. The 1.5B model achieves the deepest loss reduction, consistent with its higher capacity enabling more directional distribution shift.}
\label{tab:loss}
\small
\begin{tabular}{@{}lcccc@{}}
\toprule
Config & Gen0 loss & Gen2 loss & Final loss & Drop (\%) \\
\midrule
C8 (0.6B Full-FT K=16) & 0.206 & 0.129 & 0.097 & 53 \\
1.5B LoRA K=4 curated & 0.195 & 0.102 & 0.065 & 67 \\
\bottomrule
\end{tabular}
\end{table}

\subsection{Measurement Blind Spot}
\label{sec:blindspot}

The simultaneous stability of Vendi Score and degradation of format metrics at 0.6B reveals a measurement blind spot. Across all 0.6B configurations, Vendi Score (computed on sentence embeddings) remains within $\pm$1\% of baseline while opening entropy drops by 0.75--1.13 bits and near-duplicate rates increase by 30--37 percentage points. Embedding-based diversity metrics are insensitive to format-level homogenization because semantically distinct responses that share identical opening patterns project to different regions of embedding space. This motivates the use of complementary surface-level format metrics (opening entropy, near-duplicate rate, trigram repetition) alongside embedding-based measures for monitoring iterative training pipelines.

% ============================================================
% 5. DISCUSSION
% ============================================================
\section{Discussion}

\paragraph{The scale threshold interpretation.}
The most striking result is the qualitative difference between 0.6B (uniformly stable) and 1.5B (monotonically collapsing) under identical curation protocols. We interpret this as evidence for a scale-dependent phase transition in curation risk. At 0.6B, the model's limited capacity (further constrained by LoRA's low-rank subspace) prevents it from executing directional distribution shift in response to curation pressure. The embedding space is too constrained for the curation signal to carve out a narrowing mode. At 1.5B, the model has sufficient capacity to translate curation pressure into sustained distributional narrowing.

\paragraph{An alternative interpretation: incapacity, not resilience.}
The 0.6B stability admits a less favorable reading. If the model is simply too small to drift in any direction, then the absence of collapse is not evidence of resilience but of incapacity. The $K{=}16$ stress test (C8), where even 53\% loss reduction and 6.25\% data retention fail to produce any diversity change, strengthens this alternative. A truly resilient system would show some perturbation under extreme pressure and recover; a system that shows zero response may simply lack the capacity to respond. We cannot definitively distinguish these interpretations with our current data. We argue that both are informative for practitioners: whether the sub-1B regime is ``safe because resilient'' or ``safe because incapable,'' the practical consequence is the same.

\paragraph{The missing 1.5B control.}
Our 1.5B random baseline was terminated three times during Gen1 training due to CUDA out-of-memory errors. Random selection retains the full length distribution of generated responses, including long-tail outliers that cause activation memory spikes on our 24GB hardware. Curated selection (via ArmoRM top-1) implicitly truncates this length tail, keeping peak memory manageable. This infrastructure limitation means we cannot determine whether the 1.5B collapse is caused by curation specifically or is inherent to iterative self-training at this scale. Two interpretations remain equally viable:

\textit{Hypothesis A (curation-caused):} The reward signal at 1.5B is strong enough to drive directional mode collapse that random selection would avoid. This would mean Claim 2 (format collapse is inherent, not curation-caused) breaks at scale, and curation becomes the causal agent.

\textit{Hypothesis B (scale-inherent):} Iterative self-training at 1.5B causes collapse regardless of curation, due to the model's capacity to overfit its own distribution. Random selection would show similar degradation, just without the quality improvement.

Both hypotheses carry practical implications. Under Hypothesis A, practitioners should avoid or carefully monitor reward-based curation for models above 1B parameters. Under Hypothesis B, the entire iterative self-training paradigm requires diversity safeguards at scale, regardless of curation strategy. We leave resolution of this question to future work with sufficient compute to complete the 1.5B random baseline.

\paragraph{Scope of the causal decomposition.}
Our causal decomposition (curation vs. iterative training as drivers of format collapse) is verified at 0.6B through controlled same-$K$ comparison (Section~\ref{sec:format}). We emphasize that extrapolation of this decomposition to 1.5B or larger scales is not supported by our data. The finding that ``curation does not cause format collapse'' is a 0.6B-specific result.

% ============================================================
% 6. LIMITATIONS
% ============================================================
\section{Limitations}

\textbf{Single-model primary analysis.} Our systematic eight-configuration study is conducted on a single model (Qwen3-0.6B). While the cross-scale probe on Qwen2.5-1.5B suggests scale-dependence, this involves a different model family (Qwen2.5 vs. Qwen3), confounding scale and architecture effects.

\textbf{Incomplete 1.5B control.} The 1.5B random baseline could not be completed due to hardware constraints. This is the most significant limitation: the 1.5B collapse cannot be causally attributed to curation versus scale-inherent degradation.

\textbf{Limited scale range.} We study only 0.6B and 1.5B parameters. The location and sharpness of the hypothesized phase transition boundary cannot be characterized with two data points.

\textbf{Generation depth.} The 0.6B study uses three generations; the 1.5B study uses five. While training loss plateaus suggest convergence, longer runs might reveal slow-accumulating effects invisible in our time horizon.

\textbf{Prompt set size.} We use 2,200 prompts per generation. Larger prompt sets might alter the collapse dynamics by providing more distributional support.

\textbf{$K \times$ scale interaction.} We test $K \in \{4, 8, 16\}$ at 0.6B but only $K{=}4$ at 1.5B. The interaction between selection pressure and model scale is unexplored. $K{=}8$ or $K{=}16$ at 1.5B might produce more severe collapse.

% ============================================================
% 7. CONCLUSION
% ============================================================
\section{Conclusion}

We present a systematic empirical study of diversity dynamics under iterative reward-based curation. Our findings identify a scale-dependent phase transition: curation is benign at 0.6B parameters (eight configurations, including extreme $K{=}16$ pressure) but produces clear diversity collapse at 1.5B parameters (Vendi Score $-4.23\%$ over five generations with co-occurring collapse indicators). Format homogenization is real but inherent to iterative self-training rather than caused by curation, with selection pressure $K$ amplifying its severity (verified at 0.6B). The apparent diversity suppression under LoRA is a capacity artifact that reverses under full fine-tuning. These results suggest that practitioners working with sub-1B models need not fear curation-induced semantic collapse, but models above this threshold require careful monitoring. The incomplete 1.5B control condition prevents definitive causal attribution and represents the most important direction for future work.

% ============================================================
% REFERENCES (placeholder)
% ============================================================
\bibliographystyle{plainnat}
\begin{thebibliography}{10}

\bibitem{shumailov2023}
Shumailov, I., Shumaylov, Z., Zhao, Y., Gal, Y., Papernot, N., \& Anderson, R. (2023).
The curse of recursion: Training on generated data makes models forget.
\textit{arXiv preprint arXiv:2305.17493}.

\bibitem{alemohammad2023}
Alemohammad, S., Casco-Rodriguez, J., Luzi, L., Humayun, A.I., Babaei, H., LeJeune, D., Baraniuk, R.G. (2023).
Self-consuming generative models go MAD.
\textit{arXiv preprint arXiv:2307.01850}.

\bibitem{dohmatob2024}
Dohmatob, E., Feng, Y., \& Ren, P. (2024).
Model collapse demystified: The case of regression.
\textit{arXiv preprint arXiv:2402.07712}.

\bibitem{gao2023}
Gao, L., Schulman, J., \& Hilton, J. (2023).
Scaling laws for reward model overoptimization.
\textit{Proceedings of ICML}.

\bibitem{friedman2023}
Friedman, D. \& Dieng, A. (2023).
The Vendi Score: A diversity evaluation metric for machine learning.
\textit{Transactions on Machine Learning Research}.

\bibitem{hu2022}
Hu, E.J., Shen, Y., Wallis, P., Allen-Zhu, Z., Li, Y., Wang, S., Wang, L., \& Chen, W. (2022).
LoRA: Low-rank adaptation of large language models.
\textit{Proceedings of ICLR}.

\bibitem{wang2024armorm}
Wang, H., Tunstall, L., Groeneveld, D., \& Hajishirzi, H. (2024).
ArmoRM: Interpretable and efficient reward modeling.
\textit{arXiv preprint}.

\bibitem{zhu2024selfbleu}
Zhu, Y., Lu, S., Zheng, L., et al. (2018).
Texygen: A benchmarking platform for text generation models.
\textit{Proceedings of SIGIR}.

\end{thebibliography}

\end{document}
